\section{Capítulo~\ref{sec:debugging}. Depurando a máquina}

O que o eletrodo ouve, as poucas dimensões da atividade, o desempenho das interfaces com matrizes rígidas, o aprendizado conjunto do cérebro e do decodificador e os pontos visuais produzidos por estimulação foram medidos em experimentos revisados por pares; as estimativas dos fluxos de dados são aritmética do capítulo; sobre o dispositivo da Neuralink implantado em pessoas, até onde foi possível verificar, os dados foram publicados sobretudo pela própria empresa.

{\footnotesize
\setlength{\tabcolsep}{3pt}\renewcommand{\arraystretch}{1.15}
\begin{longtable}{>{\raggedright}p{0.5cm}>{\raggedright}p{4.0cm}>{\raggedright}p{5.6cm}>{\raggedright}p{2.3cm}>{\raggedright\arraybackslash}p{1.7cm}}
\toprule
N.º & Proposição & Experimento e o que foi medido & Fonte & Status \\
\midrule\endfirsthead
\toprule
N.º & Proposição & Experimento e o que foi medido & Fonte & Status \\
\midrule\endhead
1 & O eletrodo ouve os disparos dos neurônios num raio da ordem de cem micrômetros, em geral de um a alguns; eles são distinguidos pela forma & Rato, campo CA1, registro intracelular e extracelular simultâneo: a forma do disparo extracelular reproduz a largura e a amplitude do intracelular. Estimativa: um tetrodo poderia distinguir $60$--$100$ neurônios, mas na prática se isolam cerca de seis. & \cite{Henze2000extra} & medido \\
2 & O cérebro tem $8.6\cdot10^{10}$ neurônios; a sonda ouve cerca de um centésimo de milionésimo & Contagem de núcleos numa suspensão de tecido dissolvido: $86\cdot10^9$ neurônios no cérebro humano. A fração é aritmética do capítulo. & \cite{Azevedo2009} & medido \\
3 & A atividade de centenas de neurônios do córtex motor cabe em uma ou duas dezenas de variáveis comuns & Revisão de registros no córtex motor de macacos: a atividade da população é descrita por poucas variáveis latentes, comuns a muitos neurônios. & \cite{Gallego2017} & medido \\
4 & O ruído dos vizinhos é correlacionado, e cem neurônios carregam quase tanto quanto mil (proposição~\ref{prop:pooling}) & Córtex visual de macaco: coeficiente de correlação do ruído de neurônios vizinhos em torno de $0.1$, e o ganho da média satura com uma centena de neurônios. Teoria das correlações que limitam a informação. & \cite{Zohary1994, MorenoBote2014} & medido \\
5 & Fluxo bruto de $2\cdot10^8$~bit/s contra $8\cdot10^4$~bit/s em disparos~\eqref{eq:probedata} & Aritmética do capítulo pela capacidade do fluxo de disparos~\eqref{eq:spikeentropy}. Os parâmetros de digitalização são reais: o chip da Neuralink usa $19.3$~kHz e $10$~bits por canal. & dedução no capítulo; \cite{Rieke1997, Musk2019} & teoria \\
6 & Primeiro sistema da Neuralink: os chips amplificam e digitalizam, o fluxo bruto sai por cabo, os disparos são detectados por um programa externo; detecção no chip, carcaça fechada, rádio e energia sem fio estão no dispositivo para pessoas, segundo a empresa & Artigo da empresa: o fluxo bruto completo sai por cabo USB-C, e os disparos são detectados em tempo real por um programa fora do chip; a compressão embarcada, a alimentação e a transmissão sem fio são citadas como plano para dispositivos clínicos. Para o dispositivo implantado em pessoas, só há comunicados da empresa. Que a detecção de disparos no chip é necessária é uma conclusão do capítulo a partir de~\eqref{eq:probedata}. & \cite{Musk2019}; relatórios da Neuralink, sem artigo revisado por pares & medido (artigo da empresa); em pessoas, relatório da empresa \\
7 & Até $3072$ eletrodos em $96$ fios de $32$; os fios são inseridos por um robô que desvia dos vasos & Artigo da empresa: $3072$ eletrodos em $96$ fios de $32$; o robô insere $6$ fios ($192$ eletrodos) por minuto; em ratos com matriz implantada por longo prazo, até $70\,\%$ dos eletrodos ouvem disparos. & \cite{Musk2019} & medido (artigo da empresa) \\
8 & Em pessoas, cerca de mil canais em $64$ fios; parte dos fios se deslocou; cursor da ordem de $10$~bit/s & Só comunicados da empresa. Uma busca no PubMed não encontrou artigo revisado por pares com resultados em pessoas; isso concorda com a ressalva no texto do capítulo. & relatórios da Neuralink; sem artigo revisado por pares & medido (relatório da empresa) \\
9 & Uma interface com matriz de cem eletrodos controla um cursor & Pessoa com paralisia, $96$ eletrodos no córtex motor primário: a intenção de mover o braço altera os disparos três anos após a lesão; o decodificador fornece um ``cursor neural'' (e-mail, televisão), controle de uma prótese de mão e de um braço robótico. & \cite{Hochberg2006} & medido \\
10 & Escrita com a ``mão mental'': cerca de $90$ caracteres por minuto, menos de $8$~bit/s & Pessoa com paralisia da mão, tentativas de escrever à mão, decodificador recorrente: $90$ caracteres por minuto com precisão de $94.1\,\%$ em tempo real, acima de $99\,\%$ após autocorreção. A estimativa em bits é aritmética do capítulo. & \cite{Willett2021} & medido \\
11 & Tentativas de falar são convertidas em texto a cerca de $60$ palavras por minuto & Pessoa que perdeu a fala inteligível, matrizes no córtex: $62$ palavras por minuto; $9.1\,\%$ de erros de palavra com vocabulário de $50$ palavras, $23.8\,\%$ com vocabulário de $125\,000$. & \cite{Willett2023} & medido \\
12 & A interface extrai uma pequena fração da capacidade: um neurônio dá dezenas de bits por segundo, cem dão milhares (proposição~\ref{prop:probe}) & O limite superior~\eqref{eq:spikeentropy} é teoria; a comparação com as linhas~8 e 10 é conclusão do capítulo. Que o gargalo sejam as poucas dimensões e o desconhecimento do código, e não o fio, não foi verificado por experimento direto. & \cite{Rieke1997} & teoria \\
13 & O cérebro e o decodificador aprendem um com o outro & Macacos controlam um cursor; o decodificador se ajusta em malha fechada, e os neurônios reaprendem ao mesmo tempo: a precisão aumenta, a habilidade se mantém e sofre menos com os movimentos do próprio braço. O conselho de separar as velocidades de aprendizado é uma regra de engenharia; o capítulo não traz um experimento específico. & \cite{Orsborn2014coad} & medido \\
14 & A corrente num eletrodo excita os neurônios e fios em volta sem distinguir tipo & Córtex de camundongo e gato, registro de cálcio por dois fótons: mesmo uma corrente fraca excita neurônios esparsos, espalhados a distâncias de até milímetros, por excitação direta de axônios num volume de dezenas de micrômetros de diâmetro. Ou seja, a excitação não é seletiva, mas também não é contínua. & \cite{Histed2009stim} & medido \\
15 & A estimulação do córtex visual acende pontos; até $16$ pontos ao mesmo tempo permitem reconhecer algumas letras e bordas de objetos & Pessoa cega, $96$ eletrodos no córtex visual por $6$ meses: limiar de um eletrodo de $66.8\pm36.5$~\textmu A; a estimulação simultânea de grupos (até $16$ eletrodos, o limite do estimulador) reduz o limiar e produz padrões distinguíveis; algumas letras e bordas de objetos são reconhecidas. & \cite{Fernandez2021} & medido \\
16 & A segunda linha da Neuralink é um dispositivo que envia sinais ao córtex visual & Só comunicados da empresa sobre o desenvolvimento; não foram encontrados dados revisados por pares sobre seu funcionamento. & relatórios da Neuralink & hipótese \\
17 & As concentrações dos neurotransmissores de difusão podem ser medidas por um sensor químico no tecido & Fatias de cérebro de rato, voltametria cíclica rápida com fibra de carbono: medidas a liberação e a recaptação de serotonina; a substância sai para além da sinapse. & \cite{Bunin1998} & medido \\
\bottomrule
\end{longtable}}
